\begin{theorem}[LSeries totient eq]
\label{LSeries_totient_eq}
\lean{ArithmeticFunction.LSeries_totient_eq}
% at /home/runner/work/trough/trough/library/PrimeNumberTheoremAnd/IwaniecKowalskiCh1.lean:1845.0-1883.17
\leanok
Euler totient series: $\sum_{n=1}^{\infty} \varphi(n) n^{-s} = \zeta(s-1)/\zeta(s)$ for $\Re(s) > 2$.
  \begin{verbatim}
  This is IK (1.35). The L-series converges absolutely for Re(s) > 2
  (using φ(n) ≤ n to bound terms; the original hypothesis 1 < Re(s) is incorrect).
  \end{verbatim}
\end{theorem}
\begin{proof}
\uses{LSeries_powR_eq,powR}
\leanok
The Euler totient function $\varphi(n)$ counts the positive integers up to $n$ that are relatively prime to $n$. It is a multiplicative function, and its value at prime powers is given by $\varphi(p^k) = p^k - p^{k-1}$. The Dirichlet series of $\varphi$ can be expressed as an Euler product over primes:
\[
L(\varphi, s) = \prod_{p} \left(1 + \varphi(p)p^{-s} + \varphi(p^2)p^{-2s} + \ldots\right) = \prod_{p} \left(1 - p^{-s  +1}\right)^{-1} \left(1 - p^{-s}\right) = \frac{\zeta(s-1)}{\zeta(s)}.
\]
  The L-series converges absolutely for Re(s) > 2, using φ(n) ≤ n to bound the terms.
\end{proof}
